\begin{afterword}
\summaryhead

The essay argues that when you are wrong in a large way, you should admit it all at once. Its example is Enron, whose executives, the author says, never admitted a large mistake: they patched each failure, hid problems on the balance sheet, and after the bankruptcy Jeff Skilling told Congress that Enron had been a great company.

The author says the ``Traditional Rationality'' he grew up with, learned from Heinlein and Feynman, taught people to give in to evidence eventually but not quickly. He switched to ``Bayescraft'' after making a large mistake of his own, and realized in hindsight that he had admitted it in small, grudging concessions. He could have moved faster by saying ``oops.''

Admitting a big mistake is painful but can change your life. Getting it right the first time is better, but if you are wrong, see it all at once. Many small admissions only protect your pride and set you up to repeat the error.

\responsehead

The essay preaches confession and does not confess. Its author made ``a large mistake,'' learned the lesson of the essay from it, and then writes ``well, it's a long story'' in place of the mistake. The reader is asked to accept the method on the strength of a case they are not allowed to see. The rule the essay sets, to acknowledge a large error fully rather than in palatable pieces, is the rule it breaks in its sixth paragraph.

What the reader is shown instead is Enron, and Enron is the wrong example. The essay presents its executives as people who could not bring themselves to say ``I've been stupid.'' They were people who knew and hid: the essay's own invented dialogue has them asking how to hide a problem on the balance sheet. Skilling was convicted of conspiracy and securities fraud more than a year before the essay appeared. His testimony that he thought the company ``great'' and financially strong was a denial of knowledge, not a failure of insight. A case of deliberate deception is used to illustrate self-deception. The difference matters: the remedy for fraud is not a braver inner monologue.

The rest of the support is thin. ``Traditional Rationality'' is a position the author names, summarizes as ``fuzzy verbalisms,'' and dismisses, without quoting anyone who holds it; the first item on its list, bold and falsifiable theories, is Popper's, and Popper goes uncredited. The claim that one large loss feels better than many small ones is stated as fact, without support.

Look finally at what the essay equips the reader to do. Its last full paragraph turns to ``others,'' who concede by millimeters and say ``I was right in principle'' or ``It could have worked.'' Those phrases are handed to the reader in italics, ready to be recognized in someone else's mouth. Nothing in the essay helps a reader locate their own large mistake. The reader leaves able to diagnose other people's pride.

In short: an essay on confessing large mistakes that withholds its author's, illustrates self-deception with a fraud, and teaches the reader to spot the failure in others.
\end{afterword}
